\documentclass[10pt]{article}
% Official TMLR style files (JmlrOrg/tmlr-style-file); see STYLE_PROVENANCE.json.
% Anonymous submission: keep \usepackage{tmlr}. Use [preprint] for arXiv, [accepted] for camera-ready.
\usepackage{tmlr}
\usepackage{amsmath,amssymb}
\usepackage{booktabs}
\usepackage{graphicx}
\usepackage{hyperref}
\usepackage{url}
\usepackage{placeins}
\graphicspath{{generated/}{figures/}}

\newcommand{\R}{\mathbb{R}}
\newcommand{\ip}[2]{\langle #1,#2\rangle}

\title{Proximity Is Not Specificity: What Four Painter Names Add\\in Text-to-Image Generation}

\author{\name Anonymous authors \email anonymous@example.com \\
      \addr Paper under double-blind review}

\def\month{MM}
\def\year{YYYY}
\def\openreview{\url{https://openreview.net/forum?id=XXXX}}

\begin{document}

\maketitle

\begin{abstract}
Artist-style prompting is often evaluated by how much closer images generated with an
artist's name move to that artist's works. We show that, for related painters, this
proximity gain mostly measures a change shared by every name in the prompt set, and that
this would remain true for a generator that reproduced each painter exactly. Fixed scene
descriptions were rendered without a painting instruction, with a generic oil-painting
instruction, and in the style of Monet, Sisley, Pissarro or C\'ezanne, twice each, so that
squared magnitudes can be estimated without the bias of sampling noise. In 1,008 images from
six commercial configurations, 66.7--88.4\% of the squared change that the names add beyond
the generic instruction is shared by all four names in 31 interpretable image features. A
generator that moved the same generic outputs exactly onto each painter's reference mean
would show 84.8--95.2\%. In CLIP and CSD embeddings, the shared change supplies
73.2--83.8\% and 54.2--79.7\% of the gain in similarity to the prompted painter's
prototype, close to the faithful-imitation values of 71.3--82.9\% and 48.3--75.7\%. Painter
specificity has to be read from the differences between names instead: they align with the
reference painter differences in aggregate in every configuration, but unevenly across
painter pairs and representations. Proximity, agreement with the reference differences and
recognition of the prompted name identify different best configurations, several of them
stably under scene resampling. A separate collection of 2,000 SD-Turbo images is also
majority-shared overall, but not in texture features. Beyond the prespecified
tests, these analyses are descriptive and retrospective. We recommend reporting the
between-name comparison and a faithful-imitation benchmark next to any proximity score.
\end{abstract}

\section{Introduction}\label{sec:intro}

Artist names are among the most common style controls in text-to-image prompts
\citep{wang2023diffusiondb}, and their effect is measured in several ways. Style descriptors
score generated images by similarity to an artist's works \citep{somepalli2024style};
recognition studies ask whether the prompted artist can be identified from the image
\citep{casper2023measuring,su2025prompted,moayeri2025rethinking}; and studies of style
mimicry, protection and removal ask whether a style is reproduced, blocked or erased
\citep{shan2023glaze,honig2025adversarial,gandikota2023erasing,kumari2023ablating,zhang2024unlearncanvas}.
A common readout is proximity: images generated ``in the style of Monet'' are closer to
Monet's paintings than images generated without the name.

Proximity gain is ambiguous. When a prompt set contains related painters, a name can move
every output toward the group as a whole, without separating the painters from one another.
Such a shared movement raises similarity to every painter in the group, including the
prompted one. Whether a generator reproduces what distinguishes one painter from another is
a different question.

We make the distinction measurable with a controlled design. Fixed scene descriptions are
rendered with no painting instruction, with a generic oil-painting instruction, or with the
same instruction in the style of one of four related painters: Claude Monet, Alfred Sisley,
Camille Pissarro and Paul C\'ezanne. Averaging the four named conditions splits each name's
change into a component shared by all four names and between-name differences, the
departures from that average (Figure~\ref{fig:schematic}). Only the between-name
differences can reproduce differences among painters. Two independent generations per
condition allow squared magnitudes to be estimated from products of independent repeats,
which removes the positive bias that sampling noise adds to a squared estimate. Because a
large shared component is expected for related painters, we compare every observed quantity
with two benchmarks: an exchangeable null in which the names share nothing, and a faithful
imitator whose named outputs land exactly on each painter's reference mean.

The main collection contains 1,008 images from six text-to-image configurations, four of
them OpenAI GPT Image variants, crossed with 14 scenes, six clauses and two repeats. We
measure 31 interpretable color, spatial and texture features and, for comparison, CLIP
\citep{radford2021} and CSD \citep{somepalli2024style} embeddings of the same images. A
separate, earlier collection of 2,000 SD-Turbo images with a different template and matched
seeds provides a check outside this design. We find:
\begin{itemize}
\item \textbf{Most of what the names add is shared, and a faithful imitator would share
more.} Beyond the generic instruction, 66.7--88.4\% of the squared change added by the names
is common to all four names in the 31 features, against 84.8--95.2\% for a faithful
imitator starting from the same generic outputs. The shared change points toward the
painters' reference centroid (cosine 0.57--0.86) and covers 24.8--64.6\% of the distance
along that direction (Section~\ref{sec:shared}).
\item \textbf{Proximity gain therefore mostly measures shared movement.} In CLIP and CSD,
the shared change supplies 73.2--83.8\% and 54.2--79.7\% of the gain in similarity to the
prompted painter's prototype, within a few points of the faithful-imitation values
(Section~\ref{sec:learned}).
\item \textbf{Specificity has to be read from the between-name differences.} In every
configuration they point along the reference painter differences in aggregate, but
unevenly across painter pairs and representations. Only FLUX.2 Max's error falls clearly
below that of a generator making no painter distinctions, and only when the reference works
are resampled rather than the scenes (Section~\ref{sec:agreement}).
\item \textbf{Different readouts favor different configurations.} Nano Banana 2 has the
largest CLIP proximity gain in 85.8\% of scene resamples, while GPT Image 2 has the best CLIP
recognition in 98.7\%; FLUX.2 Max has the lowest 31-feature error in 84.7\% but the least
recognizable names (Section~\ref{sec:readouts}).
\item \textbf{Texture is the least shared feature family,} at 48.6--79.8\% in the main
collection and 36.6\% in the SD-Turbo collection (Section~\ref{sec:families}).
\end{itemize}

These results concern four related painters and digital measurements of reproductions, not
perceived resemblance, and apart from the prespecified tests the analyses are retrospective
(Section~\ref{sec:inference}). The design cannot separate movement toward these painters
from a generic effect of any artist name. The contribution is the controlled measurement,
the benchmarks that make it interpretable, and what they imply for evaluation practice.

\begin{figure}[t]
\centering
\includegraphics[width=0.92\linewidth]{fig_schematic.pdf}
\caption{The decomposition, drawn schematically (positions are illustrative, not data).
Relative to the generic clause, the four named outputs move by a shared change $c$ and
differ from one another by between-name differences $e_a$. The reference means of the four
painters differ from their centroid $\bar\mu$ by $r_a$. A faithful imitator would move the
generic output all the way along $t$ and reproduce $r_a$ exactly; because the painters are
related, most of that movement would also be shared.}
\label{fig:schematic}
\end{figure}

\section{Related work}\label{sec:related}

\paragraph{Artist style similarity and recognition.}
\citet{somepalli2024style} train contrastive style descriptors (CSD) and score generated
images against averaged artist prototypes. \citet{casper2023measuring} classify Stable
Diffusion imitations of 70 artists with CLIP and find them recognizable on average.
\citet{su2025prompted} benchmark identifying the prompted artist across 110 artists, with
content-controlled name substitutions and comparisons with the image generated from the same
seed without the artist's name. \citet{moayeri2025rethinking} consider an artist's style at
risk when the artist's own works are recognizable and the style is also recognized in
generated images, which applies to about a fifth of 372 artists. \citet{frochte2026csd}
shows that raw CSD cosine fails to separate same-artist from different-artist pairs for part
of a 91-artist corpus and also evaluates recognition of prompted generations. Artist-free
comparisons are therefore established. We add a generic-style control, split what the names
add into shared and between-name parts, and benchmark both against a faithful imitator.

\paragraph{Style mimicry, protection and removal.}
Glaze \citep{shan2023glaze} perturbs artworks to disrupt fine-tuned mimicry, and
\citet{honig2025adversarial} show with a user study that such protections can be bypassed.
Concept erasure \citep{gandikota2023erasing} and concept ablation
\citep{kumari2023ablating}, which maps an artist's style onto an anchor concept such as a
generic painting, remove styles from diffusion models; UnlearnCanvas
\citep{zhang2024unlearncanvas} benchmarks such removal, and \citet{verma2025imitation}
estimate how many training images are needed before a model imitates an artist. The
generic-painting anchor of concept ablation is close to our generic control. Our results
concern prompting of pretrained services, but they suggest asking of any such readout how
much of the measured effect every artist name would produce.

\paragraph{Expert and quantitative studies of art.}
\citet{deliege2025} assess stylistic imitation with expert raters, AI-Pastiche
\citep{asperti2025pastiche} studies authenticity and prompt adherence of generated
paintings, and \citet{fu2025aiwikiart} test vision-language models on artist attribution
and detection of generated paintings. Quantitative art research describes paintings by
color, composition and image statistics \citep{kim2014,sigaki2018,lee2020} and traces
historical change in learned embeddings \citep{kim2026}. Our features follow this tradition;
we use them to compare response components, not to define style.

\paragraph{Evaluation methodology.}
Generative evaluation separates fidelity from diversity \citep{naeem2020}, and conditional
metrics distinguish variation within a condition from relations among condition means
\citep{benny2021conditional}. Our squared-magnitude estimator is the cross-validated
distance of representational similarity analysis \citep{diedrichsen2017representational}.
Image processing can change evaluation features substantially \citep{parmar2022}, which
motivates our source-quality checks.

\section{Design and measurements}\label{sec:design}

\paragraph{Painters and reference collections.}
The reference collection contains 649 digital reproductions: 297 by Monet, 106 by Sisley,
141 by Pissarro and 105 by C\'ezanne. Works were eligible if Wikidata records them as
paintings by that painter alone, a fixed lexicon classifies their titles and metadata as
outdoor places, and Wikimedia Commons holds a reproduction at least 1,024 pixels on the
short side \citep{commonsglam,vrandecic2014wikidata}. A separate 221-work development
panel of the same painters sets the feature scaling and is never compared with generated
images; works were assigned to the two panels by a fixed hash rule before this experiment.
The painters are a deliberately hard case: three Impressionists who often painted similar
motifs, and C\'ezanne, whose work departs more clearly from theirs. Their collections differ
in subject mix (for example, 191 of 297 Monet works are water scenes, against 31 of 105 for
C\'ezanne), which we address with content-matched targets (Section~\ref{sec:sensitivity}).
Painter means receive equal weight throughout.

\paragraph{Prompts.}
Fourteen fixed outdoor scene descriptions cover water, built, land and mixed compositions
(Appendix Table~\ref{tab:scenes}). Each is crossed with six clauses: none; ``Render as an
oil painting.''; and ``Render as an oil painting in the style of Claude Monet.'' (or Alfred
Sisley, Camille Pissarro, Paul Cezanne; the prompts spell C\'ezanne without the accent). The
clause precedes the scene description, and every prompt ends with ``No text or frame.''

\paragraph{Generators.}
All requests went through the OpenRouter image API with one fixed provider per
configuration and fallbacks disabled: \texttt{openai/gpt-image-1}, \texttt{gpt-image-2},
\texttt{gpt-image-2.5-flare} and \texttt{gpt-image-2.5-sunburst} (referred to as GPT Image 1,
GPT Image 2, Flare and Sunburst; the last two are listed by the gateway as GPT Image 2.5
variants) at medium quality, Nano Banana 2 (\texttt{google/gemini-3.1-flash-image}) at 1K
resolution, and FLUX.2 Max (\texttt{black-forest-labs/flux.2-max}). All return
$1024\times1024$ images. Each configuration--scene--clause cell received two separate
requests, 1,008 images in total, collected on 10 September 2026 (UTC) in a randomized order
fixed before collection. Responses do not echo model identifiers, so the labels denote
requested configurations, not verified checkpoints. Figure~\ref{fig:examples} shows one scene
under all six clauses.

\begin{figure}[t]
\centering
\includegraphics[width=\linewidth]{specificity_original_generated.pdf}
\caption{One scene description (a river bending past a gravel bank, three trees on the far
bank and a low hill under an overcast sky) under all six clauses in each configuration
(first repeat, selected by a fixed rule before inspection). The top row shows one public-domain
reference reproduction per painter for orientation (Wikimedia Commons; titles in the
panel); references are not matched to the scene.}
\label{fig:examples}
\end{figure}

\paragraph{Measurements.}
The primary representation has 31 interpretable coordinates in three families: 11 color
features (lightness, chroma, hue concentration and entropy, and color differences at three
spatial scales), 8 spatial features (spectral slope and anisotropy, gradient statistics,
orientation entropy and balance, and quadrant variation) and 12 texture features (wavelet
energies, local-binary-pattern entropies and local contrast); Appendix~\ref{app:features}
defines them. Images are converted to sRGB and resized to a 512-pixel short side. Each
coordinate is centered and scaled by the development panel's painter-weighted median and
interquartile range. We also embed every image with CLIP ViT-L/14 \citep{radford2021} and
the released CSD ViT-L checkpoint \citep{somepalli2024style}, using each encoder's native
$224\times224$ center crop and unit-normalized 768-dimensional outputs
(Appendix~\ref{app:learned}).

\section{Separating shared and between-name change}\label{sec:method}

\subsection{The split}

Let $z_{sak}$ be one configuration's feature vector for scene $s=1,\dots,S$ ($S=14$), clause
$a$ and repeat $k\in\{1,2\}$, where $a$ is f (artist-free), g (generic) or one of the painters
$1,\dots,4$. Averaging over scenes within each repeat gives $\bar z_{ak}$. Relative to the
generic clause, the four named means split into their shared change and between-name
differences,
\begin{equation}
\bar z_{ak}-\bar z_{\mathrm{g}k}=c_k+e_{ak},\qquad
c_k=\tfrac14\sum_{a=1}^4\bar z_{ak}-\bar z_{\mathrm{g}k},\qquad
e_{ak}=\bar z_{ak}-\tfrac14\sum_{a'=1}^4\bar z_{a'k},
\label{eq:split}
\end{equation}
with squared sizes
\begin{equation}
N=4\ip{c_1}{c_2},\qquad B=\sum_{a=1}^4\ip{e_{a1}}{e_{a2}}.
\label{eq:components}
\end{equation}
If each repeat equals a fixed mean plus noise that has mean zero and is independent between
repeats, these cross-repeat products have the squared means as their expectations, whereas
squaring one repeat would add the noise variance. Estimates can therefore be negative; we
keep them. The \emph{shared fraction} is $N/(N+B)$. The between-name differences do not
depend on the baseline; against the artist-free baseline the shared component is
$N_{\mathrm{free}}=G+N+I$, where $G$ is the squared generic-minus-free shift and $I$ a
signed cross term (Appendix~\ref{app:estimators}).

\subsection{Benchmarks}

Let $\mu_a$ be painter $a$'s mean reference vector, $\bar\mu$ their centroid,
$r_a=\mu_a-\bar\mu$, and $H=\sum_a\|r_a\|^2$ the \emph{reference painter spread}. $B/H$
compares the generated between-name differences with the reference differences in size. For
the shared change we use three reference points:
\begin{itemize}
\item \textbf{Exchangeable null.} If the four named shifts were independent with mean zero
and equal expected size, the shared fraction would be 25\%.
\item \textbf{Faithful imitator.} If each named mean equalled the painter's reference mean,
the shared change would be $t_k=\bar\mu-\bar z_{\mathrm{g}k}$ and the between-name
differences $r_a$, giving the fraction $N^*/(N^*+H)$ with $N^*=4\ip{t_1}{t_2}$. This value
depends on how far the generic outputs are from the paintings, which includes the gap between
photographed paintings and generated images; it is a reference point, not an attainable
target.
\item \textbf{Direction.} The corrected cosine between $c$ and $t$, and the fraction
$\lambda=\ip{c}{t}/\|t\|^2$ of the distance to the centroid covered along its direction, both
estimated with symmetrized cross-repeat products.
\end{itemize}

\subsection{Agreement with the reference differences}

To test whether the between-name differences match the reference differences, center each
scene separately, $d_{sak}=z_{sak}-\frac14\sum_{a'}z_{sa'k}$, and define
\begin{equation}
\widehat\beta=\frac1S\sum_s\frac{1}{2}\sum_{k=1}^2\frac{\sum_a\ip{d_{sak}}{r_a}}{H},\qquad
\widehat D=\frac1S\sum_s\frac1H\sum_a\ip{d_{sa1}-r_a}{d_{sa2}-r_a}.
\label{eq:beta-d}
\end{equation}
The aligned amplitude $\beta$ projects the generated differences onto the reference pattern:
$\beta=1$ matches its size along that pattern, and $\beta>0$ indicates a response in the right
direction in aggregate, without requiring it for every pair. The error $D$ also counts
differences in all other directions: its expectation is 0 for exact agreement, 1 when the
names produce no differences, and $(\kappa-1)^2$ when the reference pattern is reproduced at
$\kappa$ times its size. With $Q=(SH)^{-1}\sum_{s,a}\ip{d_{sa1}}{d_{sa2}}$, the squared size of
the generated differences relative to the reference, $D=1-2\beta+Q$; a response in the right
direction can therefore have $D>1$ if it is too large. Because every scene is compared with
the same pooled target, $D$ also penalizes differences that vary across scenes;
$D_{\mathrm{agg}}$ scores the scene-averaged differences instead
(Appendix~\ref{app:estimators}).

\subsection{Proximity and recognition in embeddings}

In an embedding, with $g_a$ and $g_{\mathrm{g}}$ the mean unit embeddings for painter $a$ and
the generic clause and $\mu_a$ the mean reference embedding, the mean gain in cosine
similarity to the prompted painter's prototype decomposes exactly as
\begin{equation}
\underbrace{\frac14\sum_a g_a^\top\mu_a-g_{\mathrm{g}}^\top\bar\mu}_{\text{proximity gain}}
=\underbrace{(\bar g-g_{\mathrm{g}})^\top\bar\mu}_{\text{shared}}
+\underbrace{\frac{H\widehat\beta}{4}}_{\text{painter-specific}},
\label{eq:prototype-gain}
\end{equation}
with $H$ and $\beta$ computed in the embedding (Appendix~\ref{app:learned}). The shared term
does not depend on which name produced which image. Its faithful-imitation value follows by
setting $g_a=\mu_a$. Recognition is a separate readout: each named image is assigned to the
reference prototype with the largest cosine similarity, and we report macro accuracy over the
four painters. In the 31 features, the analogous proximity readout, the mean reduction in
squared distance from the generated mean to each painter's reference mean, splits exactly into
a shared term $2\ip{c}{t}-\|c\|^2$ and a between-name term
$\frac12\sum_a\ip{e_a}{r_a}-\frac14\sum_a\|e_a\|^2$ (Appendix~\ref{app:estimators}).

\subsection{Inference and status of the analyses}\label{sec:inference}

Before collection, the protocol fixed an inferential family of 21 comparisons in the 31
features, the six aligned amplitudes and the 15 pairwise differences in $D$, and declared the
artist-free shared/between-name split as a descriptive diagnostic. We use paired-scene Student
intervals with a Bonferroni adjustment across the 21, which target approximate 95\%
simultaneous coverage if scene-level errors are independent; they describe these 14 scenes
and include differences between scenes. The protocol also prescribed resampling the reference
works within painter, which we report alongside. Everything else, including the
generic-baseline split, the benchmarks, the embedding analyses, recognition, the scene
bootstrap and the SD-Turbo analysis, was added afterwards, each with its rules written in a
plan before its outcomes were computed but on images whose other results were known. Some
benchmark values were first computed during internal review of an earlier draft; the analysis
plan records this. We report these analyses with scene- or block-deletion ranges and scene-bootstrap
frequencies, and attach no significance claims to them.

\section{Results}\label{sec:results}

All 1,008 images yielded valid measurements.

\subsection{Most of what the names add is shared, and a faithful imitator would share more}\label{sec:shared}

% Generated by paper/tmlr/build_assets.py; do not edit by hand.
% input reports/painter_specificity_review_v3/analysis.json sha256 79fa5ca28900e852c072e2a52d35f3cb400ee053935d857729b4ad4aa3bc542c
% input reports/painter_tmlr_diagnostics_v1/analysis.json sha256 76792fa35f90e3d8fec1a75b0eec27d69cac14ed875fce8b64096b8e2525a4d4
\begin{table}[t]
\caption{What the four painter names add beyond the generic oil-painting clause (31 features, scene-averaged, repeat-corrected). Shared fraction: $N/(N+B)$, with its range over the 14 single-scene deletions; an exchangeable null in which the names share nothing gives 25\%. Faithful imitator: the fraction a generator would show if each named mean equalled the painter's reference mean, starting from the same generic outputs. Direction and gap covered: cosine between the shared change $c$ and the vector $t$ from the generic mean to the reference centroid, and the fraction of $t$ covered along its direction. Along generic shift: share of $N$ that continues the artist-free-to-generic shift. $B/H$: squared size of the between-name differences relative to the reference painter spread.}
\label{tab:shared}
\vspace{2pt}
\centering\small\setlength{\tabcolsep}{4pt}
\begin{tabular}{@{}lrcrrrrr@{}}
\toprule
 & \multicolumn{2}{c}{Shared fraction (\%)} & Faithful & Direction & Gap covered & Along generic & Between\\
Configuration & observed & deletions & imitator (\%) & $\cos(c,t)$ & $\lambda$ (\%) & shift (\%) & $B/H$\\
\midrule
GPT Image 1 & 71.3 & [70.4, 72.3] & 95.2 & 0.85 & 44.1 & 0.1 & 2.12\\
GPT Image 2 & 70.9 & [70.0, 72.1] & 89.3 & 0.66 & 50.7 & 0.9 & 2.04\\
GPT Image 2.5 Flare & 71.1 & [69.8, 72.3] & 84.8 & 0.57 & 54.3 & 6.6 & 2.03\\
GPT Image 2.5 Sunburst & 66.7 & [65.9, 67.7] & 86.9 & 0.72 & 55.6 & 22.0 & 1.98\\
Nano Banana 2 & 69.1 & [66.8, 70.9] & 92.5 & 0.75 & 24.8 & 55.5 & 0.61\\
FLUX.2 Max & 88.4 & [86.1, 90.2] & 90.4 & 0.86 & 64.6 & 52.8 & 0.70\\
\bottomrule
\end{tabular}
\end{table}


Beyond the generic clause, 66.7--88.4\% of the squared change that the names add is shared
by all four names (Table~\ref{tab:shared}); deleting any one scene keeps it within
65.9--90.2\%, far above the exchangeable null of 25\%. This is not by itself a sign of weak
painter specificity. A faithful imitator starting from the same generic outputs would show
84.8--95.2\%, more than every configuration (Figure~\ref{fig:benchmark}): the generic outputs
are much farther from the paintings than the four painters are from each other, so most of
any faithful movement toward them would be shared. The observed shared change points toward
the painters' reference centroid, with corrected cosine 0.57--0.86, and covers 24.8--64.6\%
of the distance along that direction.

Part of the shared change continues the generic instruction rather than moving toward these
painters. For Nano Banana 2 and FLUX.2 Max, 55.5\% and 52.8\% of $N$ lies along the
artist-free-to-generic shift, against 0.1--22.0\% for the four GPT configurations. In those
two configurations a name partly acts as a stronger painting instruction.

The between-name differences are on the scale of the reference differences: $B$ is 0.6--2.1
times the reference painter spread ($H=5.915$ in standardized units) in squared size, or 0.78--1.45 times in root-mean-square
terms. The shared change is 1.17--2.31 times the spread in root-mean-square terms. A large
shared fraction therefore does not mean that painter differences are missing: GPT Image 2's
between-name differences are about twice the reference spread in squared size, and 70.9\% of
its added change is still shared. In the 31 features, the shared change supplies all of the
gain in centroid proximity for four configurations, whose between-name differences move the
named means away from the painters' individual reference means on balance (Appendix
Table~\ref{tab:robustness}).

\begin{figure}[t]
\centering
\includegraphics[width=\linewidth]{fig_benchmark.pdf}
\caption{Shared fraction of the squared change added by the four names beyond the generic
clause (filled) and the value a faithful imitator would show from the same generic outputs
(open), in the 31 features and in each embedding. The dashed line marks the exchangeable null
of 25\%. The observed fraction is below the faithful value in every configuration and
representation.}
\label{fig:benchmark}
\end{figure}

\subsection{Proximity gain mostly measures the shared change}\label{sec:learned}

% Generated by paper/tmlr/build_assets.py; do not edit by hand.
% input reports/painter_tmlr_diagnostics_v1/analysis.json sha256 76792fa35f90e3d8fec1a75b0eec27d69cac14ed875fce8b64096b8e2525a4d4
\begin{table}[t]
\caption{The same comparison in CLIP and CSD embeddings of the same images. Shared part of proximity gain: the shared term's share of the named-minus-generic gain in cosine similarity to the prompted painter's reference prototype (Equation~\ref{eq:prototype-gain}); faithful: the same share if every named mean equalled its prototype. Shared fraction of squared change: $N/(N+B)$ computed in the embedding, with its faithful-imitation value.}
\label{tab:learned}
\vspace{2pt}
\centering\small
\begin{tabular}{@{}lrrrrrrrr@{}}
\toprule
 & \multicolumn{4}{c}{Shared part of proximity gain (\%)} & \multicolumn{4}{c}{Shared fraction of squared change (\%)}\\
\cmidrule(lr){2-5}\cmidrule(l){6-9}
 & \multicolumn{2}{c}{CLIP} & \multicolumn{2}{c}{CSD} & \multicolumn{2}{c}{CLIP} & \multicolumn{2}{c}{CSD}\\
\cmidrule(lr){2-3}\cmidrule(lr){4-5}\cmidrule(lr){6-7}\cmidrule(l){8-9}
Configuration & obs. & faithful & obs. & faithful & obs. & faithful & obs. & faithful\\
\midrule
GPT Image 1 & 74.5 & 71.3 & 72.7 & 70.4 & 76.1 & 86.1 & 76.5 & 89.2\\
GPT Image 2 & 73.2 & 76.0 & 54.2 & 48.3 & 76.6 & 87.1 & 68.8 & 82.4\\
GPT Image 2.5 Flare & 78.8 & 77.1 & 65.0 & 62.7 & 81.7 & 88.7 & 74.6 & 86.4\\
GPT Image 2.5 Sunburst & 77.7 & 77.6 & 67.4 & 65.4 & 77.6 & 88.9 & 69.9 & 86.9\\
Nano Banana 2 & 83.8 & 82.9 & 79.7 & 75.7 & 82.3 & 91.2 & 77.4 & 89.7\\
FLUX.2 Max & 81.7 & 80.7 & 77.6 & 71.0 & 80.8 & 88.5 & 75.1 & 87.7\\
\bottomrule
\end{tabular}
\end{table}


In CLIP and CSD embeddings of the same images, the shared term of
Equation~\ref{eq:prototype-gain} supplies 73.2--83.8\% (CLIP) and 54.2--79.7\% (CSD) of the
named-minus-generic gain in similarity to the prompted painter's prototype
(Table~\ref{tab:learned}). A faithful imitator would obtain 71.3--82.9\% and 48.3--75.7\% of
its gain from the same term, so the observed shares are only a few points higher. The
squared-change shares in the embeddings, 76.1--82.3\% (CLIP) and 68.8--77.4\% (CSD), are
likewise below their faithful values of 86.1--91.2\% and 82.4--89.7\%. Using the audited
source regions or the development panel as the reference target changes the shared part of
the gain by at most 1.6 percentage points. Proximity gain is thus dominated by movement that
does not depend on which name was used, for these generators and for a faithful imitator
alike; it cannot by itself show whether a model distinguishes the painters.

\subsection{Between-name differences align in aggregate, unevenly}\label{sec:agreement}

% Generated by paper/tmlr/build_assets.py; do not edit by hand.
% input reports/painter_specificity_v2/psv2-20260911/models.csv sha256 eed121f77626d8a37235b5bb19e370e3fbd7b6bae44b559d1b866e27e9cac41c
% input reports/painter_specificity_review_v1/analysis.json sha256 2fabb78e15dfcc7b027a73170f6ad0c87af0196c754b4e6f45275269c2c961e2
% input data/manifests/painter_specificity_v2/psv2-20260911/analysis.json sha256 c1d21b16f840a10fd253c1d558743a09671acf3d4e48601ddcc96625a5a90be0
\begin{table}[t]
\caption{Agreement of the between-name differences with the reference painter differences (31 features). $\beta=1$ matches the reference pattern's size along its direction; $Q$ is the squared size of the generated differences relative to the reference ($Q=1$ matches); $D=1-2\beta+Q$ is 0 for exact agreement and 1 when names produce no differences. Scenes: prespecified 95\% simultaneous intervals over 21 endpoints from paired-scene resampling. References: prespecified 95\% intervals from resampling reference works within painter, with generated vectors fixed.}
\label{tab:agreement}
\vspace{2pt}
\centering\small
\begin{tabular}{@{}lrccrrc@{}}
\toprule
 & \multicolumn{3}{c}{Aligned amplitude $\beta$} & & \multicolumn{2}{c}{Error $D$}\\
\cmidrule(lr){2-4}\cmidrule(l){6-7}
Configuration & estimate & scenes & references & $Q$ & estimate & references\\
\midrule
GPT Image 1 & 0.938 & [0.709, 1.168] & [0.767, 0.989] & 2.449 & 1.572 & [1.260, 1.892]\\
GPT Image 2 & 0.999 & [0.893, 1.104] & [0.804, 1.067] & 2.223 & 1.226 & [0.976, 1.507]\\
GPT Image 2.5 Flare & 0.772 & [0.680, 0.864] & [0.568, 0.878] & 2.281 & 1.738 & [1.407, 2.032]\\
GPT Image 2.5 Sunburst & 0.824 & [0.680, 0.968] & [0.620, 0.929] & 2.289 & 1.641 & [1.318, 1.925]\\
Nano Banana 2 & 0.442 & [0.256, 0.629] & [0.342, 0.499] & 0.994 & 1.109 & [0.961, 1.266]\\
FLUX.2 Max & 0.470 & [0.239, 0.701] & [0.369, 0.509] & 0.741 & 0.801 & [0.693, 0.945]\\
\bottomrule
\end{tabular}
\end{table}


Table~\ref{tab:agreement} compares the between-name differences with the reference
differences. Every configuration responds in the reference direction: all six aligned
amplitudes are 0.442--0.999, with prespecified scene intervals above zero and
reference-resampling lower bounds of at least 0.342. Direction and error nevertheless
diverge. GPT Image 2 matches the reference amplitude almost exactly, $\beta=0.999$ [0.893,
1.104], but its differences are more than twice the reference size, $Q=2.223$, so its error
is $D=1.226$. FLUX.2 Max has a smaller aligned response, $\beta=0.470$, smaller differences,
$Q=0.741$, and the lowest error, $D=0.801$. Its unadjusted scene interval [0.586, 1.016]
includes the no-difference value of 1, although its reference-resampling interval [0.693,
0.945] does not. Of the 15 prespecified pairwise error differences, only two are resolved:
Flare and Sunburst have higher error than FLUX.2 Max (Appendix Table~\ref{tab:pairs}).

Zero error is not attainable either. Genuine reference paintings sampled the same way, two
held-out works per painter for each of 14 pseudo-scenes and compared with a target estimated
from the other half of each collection, score $D=0.234$ on average, with a 95\% range of
[$-$0.758, 1.380] across random splits; drawing the works within content classes raises the
average to 0.753 (Appendix~\ref{app:hand}). These values reflect differences in content and
sampling between individual paintings.

The aligned response is uneven across painters. In the 31 features, Monet--Sisley alignment
ranges from $-$0.249 to 0.129 in five configurations, against 0.402 for GPT Image 2, and
C\'ezanne carries much of the aggregate response: omitting him lowers FLUX.2 Max's $\beta$
from 0.470 to 0.096 and GPT Image 2's from 0.999 to 0.651 (Appendix Table~\ref{tab:coverage}).
In both embeddings, by contrast, Monet--Sisley alignment is positive in every configuration,
0.299--0.532 in CLIP and 0.323--0.587 in CSD. The weak Monet--Sisley response in the 31
features therefore does not transfer to the embeddings; the embeddings may also respond to
content or naming cues correlated with each painter, and CSD was trained with these painters
among its style tags.

The aggregation matters as well. Scoring the scene-averaged differences gives Nano Banana 2 a
lower error than FLUX.2 Max, 0.723 against 0.761, because Nano Banana 2's differences vary
more across scenes. Shrinking each configuration's differences by one scalar fitted on the
other 13 scenes gives GPT Image 2 the lowest held-out error, 0.554 against 0.714 for FLUX.2
Max (Appendix Table~\ref{tab:calibration}). The first favors consistent differences across
scenes; the second asks how well the direction of the differences matches once their size is
corrected.

\subsection{Different readouts favor different configurations}\label{sec:readouts}

% Generated by paper/tmlr/build_assets.py; do not edit by hand.
% input reports/painter_learned_audit_v1/analysis.json sha256 06c77de4be133afc203dda122b2816d4f866e1389999f18bc94b3d9c668f28f7
% input reports/painter_specificity_v2/psv2-20260911/models.csv sha256 eed121f77626d8a37235b5bb19e370e3fbd7b6bae44b559d1b866e27e9cac41c
\begin{table}[t]
\caption{Three readouts of the same images. Proximity: named-minus-generic gain in cosine similarity to the prompted painter's reference prototype. Recognition: macro accuracy of assigning each of the 112 named images to its prompted painter by the nearest reference prototype (chance 25\%). Agreement: error $D$ of the between-name differences in each representation (lower is better). Bold marks the best displayed value in each column; Table~\ref{tab:stability} shows how often each configuration is best under scene resampling. Values are not comparable across representations.}
\label{tab:readouts}
\vspace{2pt}
\centering\small
\begin{tabular}{@{}lrrrrrrr@{}}
\toprule
 & \multicolumn{2}{c}{Proximity gain} & \multicolumn{2}{c}{Recognition (\%)} & \multicolumn{3}{c}{Agreement error $D$}\\
\cmidrule(lr){2-3}\cmidrule(lr){4-5}\cmidrule(l){6-8}
Configuration & CLIP & CSD & CLIP & CSD & 31 features & CLIP & CSD\\
\midrule
GPT Image 1 & 0.072 & \textbf{0.217} & 60.7 & 72.3 & 1.572 & \textbf{0.847} & 0.735\\
GPT Image 2 & 0.100 & 0.166 & \textbf{68.8} & \textbf{75.9} & 1.226 & 1.061 & 0.741\\
GPT Image 2.5 Flare & 0.091 & 0.187 & 59.8 & 62.5 & 1.738 & 1.057 & 0.819\\
GPT Image 2.5 Sunburst & 0.101 & 0.214 & 58.9 & 61.6 & 1.641 & 1.136 & 0.944\\
Nano Banana 2 & \textbf{0.112} & 0.210 & 51.8 & 50.9 & 1.109 & 1.078 & 0.999\\
FLUX.2 Max & 0.083 & \textbf{0.217} & 41.1 & 39.3 & \textbf{0.801} & 1.058 & \textbf{0.714}\\
\bottomrule
\end{tabular}
\end{table}

% Generated by paper/tmlr/build_assets.py; do not edit by hand.
% input reports/painter_tmlr_diagnostics_v1/analysis.json sha256 76792fa35f90e3d8fec1a75b0eec27d69cac14ed875fce8b64096b8e2525a4d4
\begin{table}[t]
\caption{Which configuration is best on each readout when the 14 scenes are resampled with replacement (5,000 paired resamples; the same scenes for every configuration). The percentages describe dependence on the authored scenes, not confidence for new scenes.}
\label{tab:stability}
\vspace{2pt}
\centering\small
\begin{tabular}{@{}llrlr@{}}
\toprule
Readout & Most often best & \% & Next & \%\\
\midrule
Proximity gain, CLIP & Nano Banana 2 & 85.8 & GPT Image 2.5 Sunburst & 9.5\\
Proximity gain, CSD & FLUX.2 Max & 33.6 & GPT Image 1 & 26.2\\
Recognition, CLIP & GPT Image 2 & 98.7 & GPT Image 1 & 0.8\\
Recognition, CSD & GPT Image 2 & 80.4 & GPT Image 1 & 19.6\\
Error $D$, 31 features & FLUX.2 Max & 84.7 & Nano Banana 2 & 15.3\\
Error $D$, CLIP & GPT Image 1 & 99.3 & GPT Image 2 & 0.6\\
Error $D$, CSD & FLUX.2 Max & 47.2 & GPT Image 1 & 26.8\\
\bottomrule
\end{tabular}
\end{table}


Table~\ref{tab:readouts} places three readouts of the same images side by side, and
Table~\ref{tab:stability} shows how often each configuration is best when the 14 scenes are
resampled. Some differences between readouts are stable. Nano Banana 2 has the largest CLIP
proximity gain, 0.112, in 85.8\% of resamples, yet its prompted names are recognized in 51.8\%
of images, while GPT Image 2 has the best CLIP recognition, 68.8\%, in 98.7\% of resamples.
FLUX.2 Max has the lowest 31-feature error in 84.7\% of resamples but the least recognizable
names, 41.1\% in CLIP and 39.3\% in CSD, below Nano Banana 2 in 99.9\% of resamples for both
encoders. GPT Image 1 has the lowest CLIP error, 0.847, in 99.3\%. Other orderings are not
stable: the configuration with the largest CSD gain is best in at most 33.6\% of resamples and
the one with the lowest CSD error in 47.2\%.

These differences have a common explanation, which follows from Equation~\ref{eq:prototype-gain}
for proximity and is an interpretation for the other two readouts. Proximity is dominated by a
term that carries no information about which name was used; recognition needs large,
separable between-name differences; and $D$ rewards differences of the right size and
direction, so it favors small differences over large misdirected ones.

A shared offset can also change recognition without changing any painter difference. Adding
one fixed translation that aligns the mean of the generated named images with the mean of the
reference prototypes, fitted on the other 13 scenes without painter labels, leaves all centered
differences unchanged but changes recognition by $-$5.4 to +12.5 percentage points in CLIP and
+0.9 to +26.8 in CSD, with means of +2.7 and +10.0 (Appendix~\ref{app:transfer}). For FLUX.2
Max in CSD it corrects 39 and harms 9 of 112 decisions. Prototypes fitted to labeled generated
images outperform both reference-based rules in all 12 configuration--encoder combinations,
so the prompted names are more separable within each generator's own outputs than with respect
to the historical references.

\subsection{Texture is the least shared feature family}\label{sec:families}

% Generated by paper/tmlr/build_assets.py; do not edit by hand.
% input reports/painter_tmlr_diagnostics_v1/analysis.json sha256 76792fa35f90e3d8fec1a75b0eec27d69cac14ed875fce8b64096b8e2525a4d4
% input reports/painter_cross_cohort_v1/analysis.json sha256 9488dfafa5779b8b065cc756804e67ad64c2841a5c62c2f9b0105f8fef13973e
\begin{table}[t]
\caption{Shared fraction $N/(N+B)$ (\%) of what the names add beyond the generic clause, by feature family (each with its own coordinates), under two alternative weightings of all 31 features, and its range when any one feature is deleted. The SD-Turbo row uses its own baseline prompt, which already requests an oil painting.}
\label{tab:families}
\vspace{2pt}
\centering\small
\begin{tabular}{@{}lrrrrrrc@{}}
\toprule
 & \multicolumn{4}{c}{Feature family} & \multicolumn{2}{c}{Weighting} & Leave one\\
\cmidrule(lr){2-5}\cmidrule(lr){6-7}
Configuration & all 31 & color & spatial & texture & equal family & covariance & feature out\\
\midrule
GPT Image 1 & 71.3 & 83.8 & 77.7 & 48.6 & 72.4 & 72.2 & [69.4, 74.6]\\
GPT Image 2 & 70.9 & 74.3 & 59.6 & 73.8 & 70.0 & 71.2 & [69.0, 73.5]\\
GPT Image 2.5 Flare & 71.1 & 81.5 & 65.5 & 60.1 & 70.8 & 76.4 & [68.4, 73.9]\\
GPT Image 2.5 Sunburst & 66.7 & 66.7 & 61.5 & 70.3 & 66.1 & 71.3 & [64.8, 70.1]\\
Nano Banana 2 & 69.1 & 56.8 & 69.8 & 79.8 & 68.9 & 65.8 & [64.8, 71.6]\\
FLUX.2 Max & 88.4 & 82.6 & 95.1 & 76.5 & 89.8 & 86.9 & [85.3, 89.6]\\
\midrule
SD-Turbo (separate collection) & 64.2 & 57.5 & 84.9 & 36.6 & -- & -- & --\\
\bottomrule
\end{tabular}
\end{table}


The shared fraction depends on the features (Table~\ref{tab:families}). In the main
collection it is 56.8--83.8\% in color and 59.6--95.1\% in spatial features, and lowest in
texture, 48.6--79.8\%, where GPT Image 1's 48.6\% is the only minority. Reweighting the 31
features gives 66.1--89.8\% (equal family weights) and 65.8--86.9\% (development covariance),
and deleting any single feature keeps it within 64.8--89.6\%.

We applied the same decomposition, with rules fixed before computing it, to all 2,000 retained
images of an earlier SD-Turbo collection: 16 short, class-level landscape descriptions (for
example, ``a riverbank landscape, with water organizing the composition''), 25 repeat blocks
and five conditions, with the baseline prompt ``An oil painting on canvas of \ldots'' and the
names inserted as ``by [painter]''. Within a scene and block, all five conditions share a
seed, so we form products only between different blocks (Appendix~\ref{app:sdturbo}). Because
the baseline already requests an oil painting, its shared fraction is comparable to
$N/(N+B)$. Across the 31 features it is 64.2\% (63.9--64.6\% across block deletions), with
57.5\% in color and 84.9\% in spatial features but 36.6\% in texture (36.2--36.9\%); computed
within scenes, it is 72.4\% overall and 48.5\% for texture. SD-Turbo's aligned amplitude is
0.618, its scene-averaged error 0.917 and its scene-wise error 1.637, and its Monet--Sisley
alignment is 0.471. The collection reuses the same references and was analyzed before in
other ways, so it is a retrospective check, not an independent replication.

\subsection{Sensitivity}\label{sec:sensitivity}

The shared fraction uses only generated images and the fixed feature scaling, so it does not
depend on the reference collection, although its benchmarks and $B/H$ do. The comparisons with
the reference differences were checked against several alternatives (Appendix
Table~\ref{tab:sensitivity}). AI assistants audited all 870 reference and development
reproductions and cropped peripheral artifacts such as frames and calibration strips from 131
of them; the crops were not verified by a human. With the cropped regions, FLUX.2 Max's error
becomes 0.832, and 0.872 after also refitting the feature scaling; all six aligned amplitudes
remain above zero, the Sunburst--FLUX.2 Max difference is no longer resolved, and the GPT
Image 1--FLUX.2 Max difference becomes resolved. FLUX.2 Max keeps the lowest error estimate
under every single-scene deletion, square measurement windows, content-matched reference
targets and both alternative feature weightings. Correcting the reference spread $H$ for its
finite-sample bias (6.7\%) rescales the aligned amplitudes to 0.474--1.071.

\section{Discussion}\label{sec:discussion}

\paragraph{What the shared change is.}
Two things contribute to it. The first is movement toward the four painters' common
appearance: the shared change points toward their reference centroid in the 31 features and
increases similarity to their average prototype in both embeddings. For related painters this
is expected, and a faithful imitator would show even more of it. The second is a
strengthening of the painting instruction itself, which accounts for about half of the shared
change for Nano Banana 2 and FLUX.2 Max. Our design cannot tell either from a generic effect
that any artist name would have; a prompt set with stylistically distant painters or a clause
naming the group (``Impressionist'') would. Generator internals, training data and guidance
could all contribute; the intervention identifies differences between prompt conditions, not
their causes.

\paragraph{Recommendations for artist-style evaluation.}
First, report the between-name comparison: whether the differences between names match the
differences between the artists, per pair as well as in aggregate. This is the readout of
specificity; proximity gain is not. Second, give proximity and shared fractions a reference
point, such as the faithful-imitation value computed from the same generic outputs, which only
needs a generic-style control and the reference means. Third, state whether errors are scored
per scene or after averaging, and whether contrasts are rescaled, since each choice changed the
ordering here. Fourth, treat recognition as its own readout, report how stable rankings are
under scene resampling, and note that recognition responds to shared offsets that leave every
painter difference unchanged.

\section{Limitations}\label{sec:limitations}

\paragraph{Scope.} The main experiment uses four related painters, one clause template, 14
authored scenes and two repeats per cell, and four of the six configurations are OpenAI
variants. The findings describe these painters and configurations, not artists in general. The
faithful-imitation benchmark depends on how far generic outputs are from photographed
paintings, which mixes style with differences in content and image formation.

\paragraph{Targets.} Our targets are digital reproductions measured by 31 features and two
related encoders. None of these establishes perceived painter resemblance, and we report no
human evaluation. The reference collections mix the paintings' properties with subject choices
and reproduction processes; the source audit addresses visible peripheral regions only and was
not verified by a human.

\paragraph{Services and inference.} The six configurations are closed services whose
checkpoints we cannot verify, and we cannot tell whether they rewrite prompts. Two repeats
cannot establish that repeated requests are independent. If repeats shared an error component,
the cross-repeat estimates of $D$ would be biased upward; under an assumed common correlation
of 0.251 the Nano Banana 2 and FLUX.2 Max error estimates would swap order
(Appendix~\ref{app:provenance}). Apart from the prespecified family, the analyses are
retrospective analyses of images whose other results were known, and the SD-Turbo collection
shares the historical references. The generated images are not public, so features cannot yet
be re-extracted from pixels by others.

\section*{Broader impact statement}
This work concerns how to interpret measurements of artist-name prompting. It may help
evaluators avoid treating proximity to an artist's collection as evidence that a model
reproduces that artist's distinguishing characteristics, which matters for debates about style
imitation. Conversely, our measurements could be misread as a quality ranking of commercial
services or as evidence about attribution, authorship or legal questions, which they are not.
We study four historical painters whose works are in the public domain; the reference
reproductions carry public-domain, CC0 or CC BY/BY-SA records, and we do not redistribute
full-resolution reference images. No living artist or human participant is involved.

\section*{Reproducibility statement}
Section~\ref{sec:design} and the appendices give the prompts, request configurations, features,
estimators and analysis rules. The supplementary archive contains the per-image features and
embeddings, the request records, every analysis output used here, and code that regenerates each
table and figure from those outputs and checks every number quoted in the text. The generated
images are not included because of the archive's size limit. AI assistants were used for code,
the reference-source audit described in Section~\ref{sec:sensitivity}, and editing; the authors
checked all reported results.

\bibliography{references}
\bibliographystyle{tmlr}

\clearpage
\appendix
\section{Prompts, requests and provenance}\label{app:provenance}

% Generated by paper/tmlr/build_assets.py; do not edit by hand.
% input data/manifests/painter_specificity_v2/psv2-20260911/requests.jsonl sha256 0aa60f491be08dbd48f1f018e225b22bcae2aacd724abf3371a6312d59acc679
\begin{table}[t]
\caption{The 14 scene descriptions, with their fixed content class. Each prompt is the clause, the description and ``No text or frame.''}
\label{tab:scenes}
\vspace{2pt}
\centering\small
\begin{tabular}{@{}rlp{0.78\linewidth}@{}}
\toprule
\# & Class & Scene description\\
\midrule
1 & water & A broad river bends past a gravel bank. Three tall trees stand on the far bank, with a low hill behind them under a pale overcast sky.\\
2 & water & A small pond lies between reeds and a grassy slope. A wooden landing extends from the left bank, with scattered clouds reflected on the water.\\
3 & water & A narrow canal passes beneath a low stone bridge. Garden walls line one side and leafy trees line the other in soft afternoon light.\\
4 & water & A quiet inlet meets a flat sandy shore. Two small boats rest near the water and distant headlands sit beneath a light grey sky.\\
5 & built & A village street curves past modest stone houses. A low wall and a single tree occupy the foreground, with hills visible beyond the roofs.\\
6 & built & A country house stands beside an orchard. A gravel path crosses the foreground and a small shed stands on the right under a clear morning sky.\\
7 & built & A square stone tower rises behind a cluster of low tiled roofs. An open field fills the foreground, with a few shrubs along a narrow footpath.\\
8 & built & A courtyard is enclosed by two simple farm buildings. A gate opens toward distant fields and a patch of sunlight falls across the bare ground.\\
9 & land & An open meadow slopes gently toward distant wooded hills. A narrow path crosses the grass and a few irregular clouds fill the upper sky.\\
10 & land & A rocky hillside carries scattered pine trees. A broad valley extends behind the foreground rocks, with distant mountains in hazy daylight.\\
11 & land & Rows of fruit trees stretch across a gently rising field. Tall grass borders the nearest row and a dark wood lies at the horizon.\\
12 & mixed & A river passes a small village at the foot of a hill. A stone wall crosses the foreground, with open grass and two trees beside the water.\\
13 & mixed & A farmhouse stands above a shallow stream. A footbridge connects two grassy banks and an orchard spreads behind the house in diffuse daylight.\\
14 & mixed & A lakeside path passes a low stone building and a cluster of trees. Water occupies the left half and a gentle ridge crosses the background.\\
\bottomrule
\end{tabular}
\end{table}


\paragraph{Scenes and prompts.}
The 14 scene descriptions (Table~\ref{tab:scenes}) are a fixed authored panel, not a sample of
possible content. Each prompt is the clause, a space, the scene description and ``No text or
frame.'' The named clauses read ``Render as an oil painting in the style of'' followed by
``Claude Monet.'', ``Alfred Sisley.'', ``Camille Pissarro.'' or ``Paul Cezanne.''; the
artist-free prompt has no clause.

\paragraph{Request configurations.}
All requests used the OpenRouter image API with a single fixed provider per configuration
(OpenAI for the four GPT Image configurations, Google AI Studio for Nano Banana 2 and Black
Forest Labs for FLUX.2 Max), fallbacks disabled, one image per request, square aspect ratio,
and no image input or seed. The OpenAI configurations requested medium quality and an opaque
background; Nano Banana 2 requested 1K resolution; FLUX.2 Max requested PNG output. The exact
payloads and a randomized request order were fixed and hashed before collection, which ran on
10 September 2026 between 15:34 and 18:21 UTC; all 1,008 requests succeeded and all images are
$1024\times1024$ pixels with distinct hashes. The responses do not echo model, quality or size
fields, so the labels identify requested configurations. The comparison concerns these
configurations, not each service's best attainable quality or equal cost per image.

\paragraph{Reference collections.}
Eligible works were recorded in Wikidata as paintings with a single creator, the painter, and
classified as outdoor places by a lexicon applied to titles and metadata that was fixed before
any feature was measured; each needed a Wikimedia Commons reproduction at least 1,024 pixels on
its short side. Works were assigned to the development and reference panels by a fixed hash
rule. By a title-based content lexicon, the reference collections comprise (water, built place,
route, open or wooded land) 191, 67, 8 and 31 Monet works; 52, 29, 9 and 16 Sisley works; 28,
77, 12 and 24 Pissarro works; and 31, 28, 12 and 34 C\'ezanne works. The Commons page cited for
the source records \citep{commonsglam} was consulted when the manuscript was written; the panels
were assembled before the experiment.

\paragraph{Repeat dependence.}
The cross-repeat estimators assume that the two repeats of a cell have independent errors, and
two repeats cannot test this. If both repeats of every cell shared an error component with a
common fraction $\rho$ of the noise power, the cross-repeat products would include that shared
component, and each configuration's error $D$ would be biased upward by an amount proportional
to its noise power. Recomputing all 15 pairwise differences under this assumption, three point
orderings reverse for $\rho<1$: Nano Banana 2 and FLUX.2 Max at $\rho=0.251$, GPT Image 1 and
GPT Image 2 at $\rho=0.246$, and GPT Image 1 and FLUX.2 Max at $\rho=0.688$. The actual
dependence is not identified, so we treat the pairwise orderings as conditional on independent
repeats. Appendix Table~\ref{tab:robustness} reports each configuration's repeat noise.

\FloatBarrier
\section{Feature definitions}\label{app:features}

\begin{table}[ht]
\caption{The 31 image features. Counts are scalar coordinates. Texture features describe the
digital image; they do not measure physical paint relief.}
\label{tab:features}
\vspace{2pt}
\centering\small
\begin{tabular}{@{}p{0.14\linewidth}p{0.42\linewidth}p{0.36\linewidth}@{}}
\toprule
Family & Measurements & Interpretation\\
\midrule
Color (11) & Lightness and chroma median and IQR; chromatic fraction; hue concentration and
entropy; color differences at three lags and their slope & Brightness, color strength and
diversity, and how quickly colors change across the image\\[3pt]
Spatial (8) & Spectral slope, residual and anisotropy; orientation entropy;
horizontal--vertical balance; gradient median and IQR; quadrant divergence & Coarse versus fine
organization, dominant directions, edge strength and regional variation in orientation\\[3pt]
Texture (12) & Four wavelet energies, their slope and curvature; three local-binary-pattern
entropies; three local coefficients of variation & Detail across scales, diversity of local
patterns and local luminance contrast\\
\bottomrule
\end{tabular}
\end{table}

Images are orientation-corrected and converted to sRGB using valid embedded color profiles;
unprofiled files are treated as sRGB. They are resized with Lanczos filtering to a 512-pixel
short side without upsampling, preserving aspect ratio. No painting-region mask is applied in
the primary analysis. The features are correlated with each other; Appendix
Table~\ref{tab:sensitivity} and Table~\ref{tab:families} report alternative weightings.

\paragraph{Color.} Features use CIELAB under D65 and the $2^\circ$ observer. Lightness and
chroma each contribute a median and an IQR. The chromatic fraction is the share of pixels with
$C^*\ge5$. Hue concentration is the length of the chroma-weighted circular mean hue; hue
entropy uses 24 equal bins, normalized by the log of the bin count, with chroma weights; both
are zero when fewer than 1\% of pixels are chromatic. At lags of 1\%, 4\% and 16\% of the short
side, the median CIEDE2000 difference \citep{sharma2005ciede} is taken over pooled horizontal
and vertical pixel pairs, and a log--log slope summarizes the three.

\paragraph{Spatial.} Features use linear-sRGB luminance. After a Tukey window (parameter 0.1),
32 logarithmic radial bins cover 4--128 cycles per short side; a Theil--Sen line gives the
spectral slope and its root-mean-square residual, and power-weighted axial concentration gives
anisotropy. Scharr derivatives give the gradient median and IQR, an 18-bin magnitude-weighted
orientation entropy and the horizontal--vertical balance. Quadrant divergence is the mean
Jensen--Shannon divergence of the four quadrants' orientation histograms from the global one.

\paragraph{Texture.} A four-level undecimated \texttt{db2} wavelet transform
\citep{mallat1989wavelet} gives four log mean-squared detail energies, their linear slope and
mean second difference. Uniform local-binary-pattern entropies \citep{ojala2002lbp} use
$(P,R)=(8,1),(16,2),(32,4)$ on quantized luminance. Three median local coefficients of
variation use windows of 1\%, 4\% and 16\% of the short side (rounded up to odd widths) with a
$10^{-6}$ floor on the mean.

\paragraph{Scaling.} Each coordinate is centered by the median and divided by the IQR of the
221-work development panel (101 Monet, 36 Sisley, 48 Pissarro and 36 C\'ezanne), with each
painter given equal total weight.

\FloatBarrier
\section{Estimators}\label{app:estimators}

\begin{table}[ht]
\caption{Estimands, their meaning and their reference values. All squared sizes use
cross-repeat products.}
\label{tab:estimands}
\vspace{2pt}
\centering\small
\begin{tabular}{@{}lp{0.52\linewidth}p{0.3\linewidth}@{}}
\toprule
Symbol & Meaning & Reference values\\
\midrule
$N$, $B$ & Squared size of the shared change beyond the generic clause, and of the
between-name differences & ---\\
$N/(N+B)$ & Shared fraction & 25\% exchangeable null; $N^*/(N^*+H)$ faithful imitator\\
$H$ & Spread of the four reference painter means, $\sum_a\|r_a\|^2$ & ---\\
$B/H$ & Size of the between-name differences relative to the reference differences & 1 for a
faithful imitator\\
$\cos(c,t)$, $\lambda$ & Direction of the shared change toward the reference centroid, and the
fraction of that distance covered & 1 and 1 for a faithful imitator\\
$\beta$ & Aligned amplitude of the between-name differences & 1 faithful; 0 no painter
distinctions\\
$Q$ & Squared size of the centered generated differences relative to $H$ & 1 faithful; 0 no
distinctions\\
$\beta/\sqrt{Q}$ & Alignment ratio (direction only) & 1 faithful\\
$D$, $D_{\mathrm{agg}}$ & Error of the scene-wise and of the scene-averaged differences & 0
faithful; 1 no distinctions\\
$V_{\mathrm{scene}}$ & Variation of the differences across scenes, $D-D_{\mathrm{agg}}$ & 0 if
constant across scenes\\
$D_{\mathrm{held}}$ & Error after rescaling the differences by a held-out scalar & ---\\
\bottomrule
\end{tabular}
\end{table}

\paragraph{Noise correction.} Write a repeat as $x_k=\theta+\eta_k$ with fixed mean $\theta$
and errors of mean zero that are independent between $k=1,2$. Then
$\mathbb E\ip{x_1}{x_2}=\|\theta\|^2$, whereas $\mathbb E\|x_1\|^2=\|\theta\|^2+\mathbb
E\|\eta_1\|^2$. Products of two different vectors are symmetrized,
$\frac12(\ip{u_1}{v_2}+\ip{u_2}{v_1})$. Centering across the four painters couples the painters
within a repeat but does not introduce dependence between repeats.

\paragraph{Exchangeable null.} If the four named shifts $h_a$ are independent with mean zero and
equal expected squared size $\sigma^2$, then $\mathbb E\,4\|\bar h\|^2=\sigma^2$ and $\mathbb
E\sum_a\|h_a\|^2=4\sigma^2$, so the expected shared fraction is $1/4$.

\paragraph{Baselines and the cross term.} Averaging scenes within repeat $k$, let $g_k$ be the
generic-minus-free shift, $c_k$ the shared change beyond the generic clause, and $g_k+c_k$ the
shared change against the artist-free baseline. Then
\begin{equation}
\underbrace{4\ip{g_1+c_1}{g_2+c_2}}_{N_{\mathrm{free}}}=\underbrace{4\ip{g_1}{g_2}}_{G}
+\underbrace{4\ip{c_1}{c_2}}_{N}
+\underbrace{4\bigl(\ip{g_1}{c_2}+\ip{c_1}{g_2}\bigr)}_{I}.
\end{equation}
The share of $N$ along the generic shift is $4\,\ip{c}{g}^2/(\|g\|^2N)$ with symmetrized
products. The within-scene variant computes the shared and between-name components in each
scene before averaging them.

\paragraph{Centroid proximity.} Write the named mean for painter $a$ as $\bar z_a=\bar
z_{\mathrm{g}}+c+e_a$ and let $t=\bar\mu-\bar z_{\mathrm{g}}$. Because $\sum_ae_a=\sum_ar_a=0$,
the mean reduction in squared distance to the painters' reference means is
\begin{equation}
\frac14\sum_a\bigl(\|\mu_a-\bar z_{\mathrm{g}}\|^2-\|\mu_a-\bar z_a\|^2\bigr)
=\underbrace{2\ip{c}{t}-\|c\|^2}_{\text{shared}}
+\underbrace{\tfrac12\textstyle\sum_a\ip{e_a}{r_a}-\tfrac14\sum_a\|e_a\|^2}_{\text{between-name}},
\end{equation}
which holds exactly with symmetrized cross-repeat products.

\paragraph{Error identities.} With $\theta_{sa}=\mathbb E[d_{sak}]$ under the repeat model,
$\mathbb E[\widehat D]=(SH)^{-1}\sum_{s,a}\|\theta_{sa}-r_a\|^2$. Writing
$\bar\theta_a=S^{-1}\sum_s\theta_{sa}$,
\begin{equation}
D=\underbrace{\frac1H\sum_a\|\bar\theta_a-r_a\|^2}_{D_{\mathrm{agg}}}
+\underbrace{\frac{1}{SH}\sum_{s,a}\|\theta_{sa}-\bar\theta_a\|^2}_{V_{\mathrm{scene}}}.
\end{equation}
Rescaling every centered generated contrast by $\kappa$ gives $D(\kappa)=1-2\kappa\beta+\kappa^2Q$;
we fit $\kappa=\max(0,\beta/Q)$ on 13 scenes and score the held-out scene.

\paragraph{Finite-sample bias of $H$.} With $n_a$ reference works and within-painter covariance
$\Sigma_a$, $\mathbb E\,\widehat H=H+\frac34\sum_a\operatorname{tr}(\Sigma_a)/n_a$; subtracting
the estimated bias lowers $H$ by 6.7\%.

\paragraph{Intervals.} For a pair of configurations, one error difference per scene gives a
paired-scene standard error; the simultaneous half-width is $t_{13,\,1-0.05/42}$ times that
standard error, and the six aligned amplitudes use the same adjustment. Their expected squared
standard error includes the variance of the fixed scene effects, so they describe these 14
scenes rather than fresh requests. In 5,000 simulated experiments per scenario with the same
design, independent Gaussian, heavy-tailed heteroskedastic and scene-interaction errors gave
simultaneous coverage of 96.0\%, 97.4\% and 99.7\%, whereas a configuration-specific state
shared by all scenes and both repeats reduced it to 0.04\%. Reference resampling draws 1,000
bootstrap samples of the reference works within painter, holding the generated vectors fixed.

\FloatBarrier
\section{Additional results for the 31 features}\label{app:hand}

% Generated by paper/tmlr/build_assets.py; do not edit by hand.
% input reports/painter_specificity_review_v3/analysis.json sha256 79fa5ca28900e852c072e2a52d35f3cb400ee053935d857729b4ad4aa3bc542c
% input reports/painter_specificity_review_v1/analysis.json sha256 2fabb78e15dfcc7b027a73170f6ad0c87af0196c754b4e6f45275269c2c961e2
% input reports/painter_tmlr_diagnostics_v1/analysis.json sha256 76792fa35f90e3d8fec1a75b0eec27d69cac14ed875fce8b64096b8e2525a4d4
\begin{table}[t]
\caption{Complete decomposition in the 31 features, in units of $H=5.915$. $G$: artist-free-to-generic shift. $N$: shared change beyond the generic clause. $I$: cross term. $N_{\mathrm{free}}=G+N+I$: shared change against the artist-free baseline. $B$: between-name differences (identical for both baselines). Within scene: $N/(N+B)$ with both components computed in each scene before averaging.}
\label{tab:decomposition-full}
\vspace{2pt}
\centering\small
\begin{tabular}{@{}lrrrrrrrr@{}}
\toprule
 & \multicolumn{5}{c}{Squared change $/H$} & \multicolumn{3}{c}{Shared fraction (\%)}\\
\cmidrule(lr){2-6}\cmidrule(l){7-9}
Configuration & $G$ & $N$ & $I$ & $N_{\mathrm{free}}$ & $B$ & vs.\ free & vs.\ generic & within scene\\
\midrule
GPT Image 1 & 5.04 & 5.26 & $-$0.29 & 10.01 & 2.12 & 82.5 & 71.3 & 71.0\\
GPT Image 2 & 6.28 & 4.98 & 1.06 & 12.32 & 2.04 & 85.8 & 70.9 & 76.2\\
GPT Image 2.5 Flare & 4.01 & 4.99 & 2.29 & 11.29 & 2.03 & 84.8 & 71.1 & 76.9\\
GPT Image 2.5 Sunburst & 3.08 & 3.97 & 3.28 & 10.32 & 1.98 & 83.9 & 66.7 & 72.2\\
Nano Banana 2 & 1.22 & 1.36 & 1.91 & 4.49 & 0.61 & 88.1 & 69.1 & 52.8\\
FLUX.2 Max & 3.78 & 5.34 & 6.53 & 15.65 & 0.70 & 95.7 & 88.4 & 93.6\\
\bottomrule
\end{tabular}
\end{table}

% Generated by paper/tmlr/build_assets.py; do not edit by hand.
% input reports/painter_tmlr_diagnostics_v1/analysis.json sha256 76792fa35f90e3d8fec1a75b0eec27d69cac14ed875fce8b64096b8e2525a4d4
\begin{table}[t]
\caption{Further diagnostics in the 31 features. Repeat noise: squared difference between the two repeats' centered contrasts, relative to $H$. $\beta$ with corrected $H$: aligned amplitude after removing the finite-sample bias of $H$ (6.7\%). $\cos(c,g)$: corrected cosine between the shared naming change and the artist-free-to-generic shift. Centroid proximity gain: mean reduction in squared distance from the generated mean to each painter's reference mean, from the generic to the named clause, split exactly into shared and between-name parts (Appendix~\ref{app:estimators}).}
\label{tab:robustness}
\vspace{2pt}
\centering\small
\begin{tabular}{@{}lrrrrrr@{}}
\toprule
 & Repeat noise & $\beta$ with & $\cos(c,g)$ & \multicolumn{3}{c}{Centroid proximity gain}\\
\cmidrule(l){5-7}
Configuration & $/H$ & corrected $H$ & & total & shared & between-name\\
\midrule
GPT Image 1 & 2.02 & 1.006 & $-$0.03 & 17.490 & 17.844 & $-$0.354\\
GPT Image 2 & 0.96 & 1.071 & 0.09 & 5.094 & 5.159 & $-$0.065\\
GPT Image 2.5 Flare & 0.42 & 0.827 & 0.26 & 0.871 & 1.585 & $-$0.715\\
GPT Image 2.5 Sunburst & 0.58 & 0.883 & 0.47 & 4.500 & 4.997 & $-$0.498\\
Nano Banana 2 & 2.60 & 0.474 & 0.74 & 7.423 & 7.013 & 0.410\\
FLUX.2 Max & 1.67 & 0.504 & 0.73 & 10.529 & 10.176 & 0.353\\
\bottomrule
\end{tabular}
\end{table}

% Generated by paper/tmlr/build_assets.py; do not edit by hand.
% input reports/painter_specificity_v2/psv2-20260911/model_contrasts.csv sha256 0d0b722750e657b83658ccffbc7cb093d7b0ae577896904c45c02ababe31665b
\begin{table}[t]
\caption{All 15 prespecified pairwise error comparisons in the 31 features. Brackets are approximate 95\% simultaneous intervals over the 21-endpoint family; a positive difference favors configuration B. $^{*}$: interval excludes zero.}
\label{tab:pairs}
\vspace{2pt}
\centering\small
\begin{tabular}{@{}llc@{}}
\toprule
Configuration A & Configuration B & $D_A-D_B$ [simultaneous interval]\\
\midrule
GPT Image 1 & GPT Image 2 & 0.346 [$-$0.400, 1.093]\\
GPT Image 1 & GPT Image 2.5 Flare & $-$0.165 [$-$1.027, 0.697]\\
GPT Image 1 & GPT Image 2.5 Sunburst & $-$0.069 [$-$1.181, 1.043]\\
GPT Image 1 & Nano Banana 2 & 0.463 [$-$0.602, 1.528]\\
GPT Image 1 & FLUX.2 Max & 0.771 [$-$0.037, 1.580]\\
GPT Image 2 & GPT Image 2.5 Flare & $-$0.512 [$-$1.150, 0.126]\\
GPT Image 2 & GPT Image 2.5 Sunburst & $-$0.415 [$-$1.181, 0.350]\\
GPT Image 2 & Nano Banana 2 & 0.116 [$-$0.922, 1.155]\\
GPT Image 2 & FLUX.2 Max & 0.425 [$-$0.161, 1.011]\\
GPT Image 2.5 Flare & GPT Image 2.5 Sunburst & 0.096 [$-$0.286, 0.478]\\
GPT Image 2.5 Flare & Nano Banana 2 & 0.628 [$-$0.283, 1.540]\\
GPT Image 2.5 Flare & FLUX.2 Max & 0.937 [0.256, 1.618]$^{*}$\\
GPT Image 2.5 Sunburst & Nano Banana 2 & 0.532 [$-$0.436, 1.500]\\
GPT Image 2.5 Sunburst & FLUX.2 Max & 0.840 [0.058, 1.623]$^{*}$\\
Nano Banana 2 & FLUX.2 Max & 0.309 [$-$0.847, 1.464]\\
\bottomrule
\end{tabular}
\end{table}

% Generated by paper/tmlr/build_assets.py; do not edit by hand.
% input reports/painter_specificity_review_v1/analysis.json sha256 2fabb78e15dfcc7b027a73170f6ad0c87af0196c754b4e6f45275269c2c961e2
\begin{table}[t]
\caption{Scene aggregation and contrast magnitude in the 31 features. $D=D_{\mathrm{agg}}+V_{\mathrm{scene}}$; $\beta/\sqrt{Q}$ is the alignment ratio (direction only). $D_{\mathrm{held}}$: error after multiplying every centered generated contrast by a nonnegative scalar fitted on the other 13 scenes; the fitted range covers the 14 folds. Rescaling changes feature vectors, not images.}
\label{tab:calibration}
\vspace{2pt}
\centering\small
\begin{tabular}{@{}lrrrrcr@{}}
\toprule
Configuration & $D$ & $D_{\mathrm{agg}}$ & $V_{\mathrm{scene}}$ & $\beta/\sqrt{Q}$ & fitted scalar & $D_{\mathrm{held}}$\\
\midrule
GPT Image 1 & 1.572 & 1.239 & 0.333 & 0.600 & 0.366--0.398 & 0.645\\
GPT Image 2 & 1.226 & 1.044 & 0.182 & 0.670 & 0.438--0.459 & \textbf{0.554}\\
GPT Image 2.5 Flare & 1.738 & 1.483 & 0.255 & 0.511 & 0.329--0.349 & 0.741\\
GPT Image 2.5 Sunburst & 1.641 & 1.337 & 0.305 & 0.545 & 0.346--0.373 & 0.706\\
Nano Banana 2 & 1.109 & 0.723 & 0.387 & 0.444 & 0.401--0.533 & 0.832\\
FLUX.2 Max & 0.801 & 0.761 & 0.040 & 0.546 & 0.598--0.694 & 0.714\\
\bottomrule
\end{tabular}
\end{table}

% Generated by paper/tmlr/build_assets.py; do not edit by hand.
% input reports/painter_specificity_review_v1/analysis.json sha256 2fabb78e15dfcc7b027a73170f6ad0c87af0196c754b4e6f45275269c2c961e2
% input reports/painter_specificity_review_v3/analysis.json sha256 79fa5ca28900e852c072e2a52d35f3cb400ee053935d857729b4ad4aa3bc542c
% input reports/painter_learned_audit_v1/analysis.json sha256 06c77de4be133afc203dda122b2816d4f866e1389999f18bc94b3d9c668f28f7
\begin{table}[t]
\caption{Which painter differences carry the aligned response. Components: aligned amplitude along each singular component of the centered reference means (66.3\%, 20.6\% and 13.0\% of their spread). Omitting a painter recenters the other three and recomputes $\beta$. Monet--Sisley $\beta$: aligned amplitude for that pair alone, in the 31 features (full frame and central-square reference windows) and in each embedding.}
\label{tab:coverage}
\vspace{2pt}
\centering\footnotesize
\begin{tabular}{@{}l*{11}{r}@{}}
\toprule
 & \multicolumn{3}{c}{Component} & \multicolumn{4}{c}{$\beta$ after omitting} & \multicolumn{4}{c}{Monet--Sisley $\beta$}\\
\cmidrule(lr){2-4}\cmidrule(lr){5-8}\cmidrule(l){9-12}
Configuration & 1st & 2nd & 3rd & Mon. & Sis. & Pis. & C\'ez. & 31 & square & CLIP & CSD\\
\midrule
GPT Image 1 & 1.248 & 0.310 & 0.356 & 1.140 & 1.072 & 0.840 & 0.389 & 0.074 & 0.038 & 0.301 & 0.446\\
GPT Image 2 & 1.156 & 0.617 & 0.800 & 1.195 & 0.972 & 0.992 & 0.651 & 0.402 & 0.408 & 0.532 & 0.587\\
GPT Image 2.5 Flare & 1.022 & 0.211 & 0.383 & 1.034 & 0.734 & 0.802 & 0.222 & $-$0.011 & 0.012 & 0.339 & 0.451\\
GPT Image 2.5 Sunburst & 1.082 & 0.250 & 0.417 & 1.116 & 0.840 & 0.762 & 0.284 & $-$0.249 & $-$0.185 & 0.385 & 0.323\\
Nano Banana 2 & 0.513 & 0.371 & 0.195 & 0.569 & 0.446 & 0.388 & 0.276 & $-$0.044 & $-$0.017 & 0.318 & 0.400\\
FLUX.2 Max & 0.666 & 0.021 & 0.183 & 0.539 & 0.511 & 0.527 & 0.096 & 0.129 & 0.132 & 0.299 & 0.414\\
\bottomrule
\end{tabular}
\end{table}

% Generated by paper/tmlr/build_assets.py; do not edit by hand.
% input reports/painter_specificity_review_v1/analysis.json sha256 2fabb78e15dfcc7b027a73170f6ad0c87af0196c754b4e6f45275269c2c961e2
% input reports/painter_reference_quality_v1/analysis.json sha256 089f636335129d6377328088ccc9fdc763488aba2a88463939322326707052b4
\begin{table}[t]
\caption{Error $D$ under alternative targets, weightings and source corrections. Class target: title-derived water, built and land reference means for the 11 matching scenes (pooled: the same 11 scenes against the pooled target). Equal family: each feature family has equal total weight. Covariance: shrunk inverse covariance of the development panel. Cropped: AI-audited painting regions with the original scaling; refit: also refitting the development scaling. Each column has its own normalizer.}
\label{tab:sensitivity}
\vspace{2pt}
\centering\small
\begin{tabular}{@{}lrrrrrrr@{}}
\toprule
 & & \multicolumn{2}{c}{11 scenes} & \multicolumn{2}{c}{Weighting} & \multicolumn{2}{c}{Source correction}\\
\cmidrule(lr){3-4}\cmidrule(lr){5-6}\cmidrule(l){7-8}
Configuration & Primary & pooled & class & equal family & covariance & cropped & refit\\
\midrule
GPT Image 1 & 1.572 & 1.557 & 1.366 & 1.551 & 2.103 & 1.561 & 1.631\\
GPT Image 2 & 1.226 & 1.099 & 1.130 & 1.243 & 1.511 & 1.171 & 1.259\\
GPT Image 2.5 Flare & 1.738 & 1.726 & 1.539 & 1.765 & 1.636 & 1.680 & 1.737\\
GPT Image 2.5 Sunburst & 1.641 & 1.633 & 1.525 & 1.680 & 1.617 & 1.536 & 1.577\\
Nano Banana 2 & 1.109 & 0.987 & 0.948 & 1.203 & 1.385 & 1.076 & 1.121\\
FLUX.2 Max & 0.801 & 0.750 & 0.786 & 0.808 & 0.928 & 0.832 & 0.872\\
\bottomrule
\end{tabular}
\end{table}


\begin{figure}[ht]
\centering
\includegraphics[width=\linewidth]{specificity_artist_pairs.pdf}
\caption{Aligned amplitude and error for each of the six painter pairs (M: Monet, S: Sisley,
P: Pissarro, C: C\'ezanne), using all 31 features and normalized by the squared reference
distance of each pair. An amplitude of 1 matches the reference, negative values point the
opposite way, and an error of 1 corresponds to no distinction. Negative errors can arise from
repeat correction.}
\label{fig:pairs}
\end{figure}

\paragraph{Full decomposition.} Table~\ref{tab:decomposition-full} gives all components. Against
the artist-free baseline, 82.5--95.7\% of the squared change is shared; this includes the
painting instruction itself, and computing both components within scenes gives 88.6--97.2\%.
Against the generic baseline the within-scene shares are 52.8--93.6\%.

\paragraph{Genuine-painting controls.} For each of 1,000 random splits, the works in each
painter's collection (or each painter and content class) are divided into two disjoint halves;
the smaller half estimates the target, and the other supplies two works per painter for each of
14 pseudo-scenes (11 with content classes), which are scored with the same error $D$. Pooled
sampling gives a mean of 0.234 with a central 95\% range of [$-$0.758, 1.380]; sampling within
content classes against the pooled target gives a mean of 0.753.

\paragraph{Painter pairs and coverage.} Table~\ref{tab:coverage} shows that the aggregate
aligned response rests mostly on the first reference component, which accounts for 66.3\% of
the reference spread, and that omitting C\'ezanne removes much of it. Stronger pair alignment
does not ensure lower pair error: GPT Image 1 has weaker Monet--Sisley alignment than GPT Image
2 but lower error for that pair (Figure~\ref{fig:pairs}). For Flare, Sunburst and Nano Banana
2, swapping the Monet and Sisley reference labels lowers the overall error, as it must when the
pair alignment is negative.

\paragraph{Scene aggregation and rescaling.} Nano Banana 2 has the largest scene variation and
FLUX.2 Max the smallest, which explains their reversal between $D$ and $D_{\mathrm{agg}}$. GPT
Image 2's differences are about twice the reference size, so shrinking them to about 0.45 of
their size removes more off-pattern difference than aligned response. Deleting any scene and
refitting the whole procedure preserves GPT Image 2's advantage after rescaling.

\paragraph{Reference sensitivity.} The content-matched target uses the title-derived water,
built and land classes for the 11 non-mixed scenes. The source audit recorded removable
frames, margins, labels and calibration strips; 43 uncertain boundaries were left unchanged. A
visual content audit of the same reproductions disagreed with the title-derived class for 230
works; replacing the title-derived with visual labels also leaves FLUX.2 Max with the lowest
error. All works and features are retained under every correction.

\FloatBarrier
\section{Learned embeddings}\label{app:learned}

\paragraph{Extraction.} We used the public checkpoints \texttt{openai/clip-vit-large-patch14}
and \texttt{tomg-group-umd/CSD-ViT-L}. CSD uses its released style projection of the CLIP visual
backbone; its release notes a discrepancy with the published benchmark numbers, so we
characterize the released checkpoint without claiming to reproduce those numbers. Every image
is orientation- and color-corrected as above, resized bicubically to a 224-pixel short side,
center-cropped to $224\times224$ and normalized with the CLIP channel statistics. Embeddings
are unit-normalized 768-dimensional vectors; we use their Euclidean and cosine geometry without
scaling or dimension selection. Backend and batch-size checks on fixed test cases agreed to
within $1.3\times10^{-6}$ per coordinate.

\paragraph{Prototype gain.} With unnormalized prototypes $\mu_a$ (means of unit reference
embeddings), $g_a^\top\mu_b$ is the average image-to-reference cosine similarity. Writing
$g_a=\bar g+d_a$ and $\mu_a=\bar\mu+r_a$,
\begin{equation}
\frac14\sum_a g_a^\top\mu_a-g_{\mathrm g}^\top\bar\mu
=(\bar g-g_{\mathrm g})^\top\bar\mu+\frac14\sum_a d_a^\top r_a
=(\bar g-g_{\mathrm g})^\top\bar\mu+\frac{H\widehat\beta}{4},
\end{equation}
because $\sum_a r_a=\sum_a d_a=0$. The first term cannot depend on which name produced which
images; permuting the painter labels of the generated images changes only $\widehat\beta$. The
faithful value sets $g_a=\mu_a$, giving the shared fraction $(\bar\mu-g_{\mathrm
g})^\top\bar\mu/[(\bar\mu-g_{\mathrm g})^\top\bar\mu+H/4]$.

\paragraph{Scope.} CSD is built on the CLIP backbone and both encoders see the same images, so
their agreement is evidence about related representations, not independent confirmation. All
four painters appear in CSD's style-tag vocabulary, and the overlap between either encoder's
training data and our references is unknown. As a check on the reference panels, the primary
prototypes classify the 221 development works with 79.8\% (CLIP) and 79.6\% (CSD) macro
accuracy.

\FloatBarrier
\section{Held-scene prompt-name recognition}\label{app:transfer}

% Generated by paper/tmlr/build_assets.py; do not edit by hand.
% input reports/painter_prototype_transfer_v1/analysis.json sha256 e1d1fbb924feb797f67e0907677511ed7c745335a4f2827e99d2f286e6be5448
\begin{table}[t]
\caption{Held-scene recognition of the prompted painter (\%, 112 named images per configuration, chance 25\%). Ref.: nearest reference prototype. Shift: the same after adding one fixed translation that aligns the generated and reference mixture means, fitted on the other 13 scenes without painter labels. Gen.: nearest generated-image prototype fitted with painter labels on the other 13 scenes. Shift and Gen.\ use more information than Ref., and Gen.\ more than Shift.}
\label{tab:transfer}
\vspace{2pt}
\centering\small
\begin{tabular}{@{}lrrrrrr@{}}
\toprule
 & \multicolumn{3}{c}{CLIP} & \multicolumn{3}{c}{CSD}\\
\cmidrule(lr){2-4}\cmidrule(l){5-7}
Configuration & Ref. & Shift & Gen. & Ref. & Shift & Gen.\\
\midrule
GPT Image 1 & 60.7 & 55.4 & 80.4 & 72.3 & 73.2 & 84.8\\
GPT Image 2 & 68.8 & 66.1 & 75.0 & 75.9 & 76.8 & 94.6\\
GPT Image 2.5 Flare & 59.8 & 59.8 & 77.7 & 62.5 & 74.1 & 92.9\\
GPT Image 2.5 Sunburst & 58.9 & 62.5 & 69.6 & 61.6 & 70.5 & 90.2\\
Nano Banana 2 & 51.8 & 59.8 & 63.4 & 50.9 & 61.6 & 75.0\\
FLUX.2 Max & 41.1 & 53.6 & 71.4 & 39.3 & 66.1 & 75.0\\
\midrule
Mean change vs.\ Ref.\ (pp) & & +2.7 & +16.1 & & +10.0 & +25.0\\
\bottomrule
\end{tabular}
\end{table}


For each configuration and held-out scene $s$, let $x$ be a unit embedding of a named image
from scene $s$ and $u_a=\mu_a/\|\mu_a\|$ the normalized reference prototypes. The reference rule
(Ref.) is $\arg\max_a x^\top u_a$. The shift rule computes $t_{-s}=\frac14\sum_a\mu_a-\bar
g_{-s}$, where $\bar g_{-s}$ is the mean of the 104 named embeddings from the other 13 scenes,
and predicts $\arg\max_a(x+t_{-s})^\top u_a$. The translation changes each class score by the
constant $t_{-s}^\top u_a$ and cancels exactly under centering across painters, so $\beta$, $Q$
and $D$ are unchanged. The generated-prototype rule (Gen.) uses the normalized mean of the 26
labeled training embeddings for each painter. Held-scene repeats never enter the fit, and the
artist-free and generic images are not used. The question was chosen after the reference-rule
confusions were known; the rules were fixed before the translated outcomes were computed. Folds
overlap and scenes are authored, so these are descriptive out-of-fold results. Across the four
combinations of source view and reference target, the mean translation gains are 2.2--4.0
points (CLIP) and 9.7--10.1 points (CSD).

\FloatBarrier
\section{The SD-Turbo collection}\label{app:sdturbo}

The collection was generated on 5 September 2026 with \texttt{stabilityai/sd-turbo} at
$512\times512$ pixels, one denoising step and guidance scale zero, for an earlier, unpublished
distance analysis by the same project. It crosses 16 scenes, 25 repeat blocks and five
conditions (the baseline and four names), 400 images per condition. Within each scene and
block, all five conditions use the same seed. Its image hashes are disjoint from the main
collection, and its 16 scene descriptions are shorter, class-level texts, different from the 14
main descriptions.

For block vectors $v_b$, $b=1,\dots,K$ with $K=25$, the estimator
$U(v)=[K(K-1)]^{-1}\sum_{b\ne b'}\ip{v_b}{v_{b'}}$ replaces the two-repeat product. Averaging
scenes within each block, let $\delta_{ba}$ be the named-minus-baseline shift, $c_b$ its mean
over names and $e_{ba}=\delta_{ba}-c_b$. The shared component is $4U(c)$ and the between-name
component is $\sum_aU(e_a)$. Because products are formed only between different blocks,
correlation among the five seed-matched conditions within a block is allowed. The reference
means and scaling are those of the main experiment. Deleting each block once gives the reported
ranges; they describe sensitivity, not uncertainty.

\FloatBarrier
\section{Evidence and replay}\label{app:replay}

Every analysis in this paper reads retained feature vectors or embeddings and writes a versioned
result file whose inputs are recorded by SHA-256. The tables and Figures~\ref{fig:schematic}
and~\ref{fig:benchmark} are generated from those result files, and a check script verifies the
input hashes, regenerates every table and figure byte for byte and confirms that each number
quoted in the text equals its recomputed value. These checks establish that the reported numbers
follow from the retained measurements; they do not re-extract features from pixels.


\end{document}
